\scanpage{294}
% archive sheet number: 295
Problème MP 2

Un des objets du présent exercice est d'étudier la convergence de certaines séries de la forme
\[
u_n/\unsure{n|m}\quad ,
\]
où $(u_n)$ est une suite bornée de nombres complexes, et d'en calculer la somme dans certains cas, notamment dans le cas où la suite $(u_n)$ admet une ``période'' $r$ (de sorte qu'on a $u_{n+r}=u_n$ pour tout $r$).

I

Soit $(c_n)_{n\quad1}$ une suite décroissante de nombres positifs tendant vers \deleted{0, et soit \unsure{(u}} 0, et soit $(u_n)_{n\quad1}$ une suite de nombres complexes. \inserted{a)} On suppose qu'il existe un nombre $M\quad0$ tel que pour tout entier $n\quad0$, on a
\[
S_n={}_{1\ i\ n}\;u_i\quad \unsure{.}\quad M\quad .
\]
\deleted{a)} Prouver que la série
\[
u_n c_n
\]
est convergente, et que pour tout entier $n\quad1$, le reste
\[
R_n={}_{i\ \,n+1}\;u_i c_i
\]
satisfait
\[
R_n\quad Mc_{n+1}\quad .
\]
(Utiliser le critère de Cauchy, en prouvant que pour $m\quad n$, on a
\[
{}_{n\ i\ m}\;u_i c_i\qquad Mc_{n+1}\quad ;
\]
pour ceci, remplacer dans le premier membre $c_i$ par $S_i-S_{i-1}$.)

b) On suppose que la suite $(u_n)$ est périodique de période $r$ ($r\quad1$), et on s'intéresse à la convergence de la série $\;u_n c_n$. Montrer que si la série $\;c_n$ est convergente, il en est de même de $\;u_n c_n$ \deleted{\unsure{q|u}}quelle que soit la suite (périodique) $(u_n)$. Montrer que si la série $\;c_n$ est divergente, alors $\;u_n c_n$ est convergente si et seulement si on a
\[
{}_{1\ i\ r}\;u_i=0\unsure{/}\quad .
\]
(Montrer que toute suite $(u_n)$ périodique de période $r$ se met de façon unique sous forme d'une somme de deux suites périodiques de période $r$
\[
u_n=u'_n+u''_n\quad ,
\]
où $(u'_n)$ est une suite constante et où $(u''_n)$ satisfait à la condition
\[
{}_{1\ i\ r}\;u''_i=0\quad .
\]
Appliquer alors le résultat de a) à cette dernière suite.)

c) On considère l'espace \inserted{vectoriel\unsure{:}complexe $E_r$} des suites périodiques \inserted{$u=$}$(u_n)$ de période $r$ satisfaisant la condition (3). Soit $n_{\unsure{o|0}}\quad1$, et soit , pour tout entier $\deleted{\unsure{i}}s$ tel que $1\quad\deleted{\unsure{i}}s\quad r-1$, $u(\deleted{\unsure{i}}s)=(u(\deleted{\unsure{i}}s)_n)_{n\quad1}$ l'unique suite $\;E_r$ telle qu'on ait
\[
\begin{aligned}
\text{pour }0\quad i\quad r-1,\quad \deleted{u(i)_{\unsure{n_o+i}}}\;u(\deleted{\unsure{i}}s)_{n_o+\unsure{i|1}}&=\begin{aligned}
1&\ \text{si }\deleted{\text{\unsure{$i$ est 0, ou $s=1$ si}}}\ i=0\\
-1&\ \text{si }i=s\\
0&\ \text{si }i\ne0,s
\end{aligned}\quad .
\end{aligned}
\]
Prouver que les $u(s)$ ($1\quad s\quad r-1$) forment une base de l'espace vectoriel
